\documentclass[11pt]{exam}
\usepackage[activeacute,spanish]{babel} % Permite el idioma espa\~nol.
\usepackage[utf8]{inputenc}
\usepackage{amsmath}
\usepackage[colorlinks]{hyperref}

\usepackage{minted} 
\usemintedstyle{emacs}
\usepackage{tcolorbox} % colores para el fondo
\definecolor{bg}{rgb}{0.95,0.95,0.95} %color de fondo

\pagestyle{headandfoot}

\begin{document}

\firstpageheadrule
%\firstpagefootrule
%\firstpagefooter{}{Pagina \thepage\ de \pages}{}
\runningheadrule
%\runningfootrule
\lhead{\bf\normalsize Computación Científica\\ Guía 06}
\rhead{\bf\normalsize Cs. Fís., Astro., Geofís. \\ 2023-1}
\chead{\bf\normalsize Depto. de Física \\ Universidad de Concepción}
%\rfoot{\thepage\ / pages}
\cfoot{ }
\lfoot{\tiny GR}
\begin{flushleft}
\vspace{0.2in}
%\hbox to \textwidth{Nombre: \enspace \hrulefill}
%Nombre : \\
\vspace{0.25cm}
\end{flushleft}
%%%%%%%%%%%%%%%%%%%%%%%%%%%%%%%%%%%%%%%%%%

%\begin{center}
%\texttt{Fecha de Entrega: Jueves 28 de Agosto. Env'ie a gfrubi@udec.cl los archivos .py que resuelven cada uno de los problemas propuestos.}
%\end{center}
\begin{questions}

\item Cree un nuevo documento \texttt{.tex} que cargue los paquetes \texttt{amsmath} y \texttt{amsfonts} 

(use \verb|\usepackage{amsmath,amsfonts}|) y que genere las siguientes expresiones matem\'aticas:
\begin{itemize}

\item La siguiente ecuaci'on numerada:
\begin{equation}\label{asinNnu}
Y_\nu(x)\approx \sqrt{\frac{2}{\pi x}}\sin\left(x-\frac{\nu\pi}{2}-\frac{\pi}{4}\right), 
\qquad x\gg\left|\nu^2-\frac{1}{4}\right|.
\end{equation}
Usando 

\begin{minted}[bgcolor=bg]{tex}
	\begin{equation}

	\end{equation}
\end{minted}
\item La siguiente ecuaci'on no numerada:
\begin{equation*}
\int_{0}^{\infty }\frac{\log (x)}{x^{2}}dx=-\left. \frac{\log (x)}{x}\right]
_{0}^{\infty }+\int_{0}^{\infty }\frac{1}{x^{2}}dx
\end{equation*}
Usando

\begin{minted}[bgcolor=bg]{tex}
	\begin{equation*}

	\end{equation*}
\end{minted}
\item La siguiente expresi'on del m'ultiples l'ineas:
\begin{align} 
(a+b)^4 &= (a+b)^2 (a+b)^2 \\
 &= (a^2+2ab+b^2) (a^2+2ab+b^2) \\
 &= a^4+4 a^3 b + 6 a^2 b^2 +4 a b^3 +b^4
\end{align}
Usando (ver páginas 47-48 de la \href{https://udec.instructure.com/courses/40179/pages/latex-pdf-presentacion?module_item_id=1465499}{presentación de clases})

\begin{minted}[bgcolor=bg]{tex}

	\begin{align} 
	 ... &= ...\\
	 ... &= ...\\
	 ... &= ...
	\end{align} 
\end{minted}

\item Una expresi'on enmarcada, usando el comando \verb|\boxed{}| de \texttt{amsmath}:
\newline
\begin{equation*}
\boxed{\int u \, dv=u\,v-\int v \, du}
\end{equation*}

\item Diferentes tipografías matem\'aticas con fuentes de \texttt{amsfonts}:
\begin{itemize}
\item 
\begin{equation}
\mathsf{C}_{ijkl}=\boldsymbol{\mathsf{C}}
\end{equation}
Usando

\begin{minted}[bgcolor=bg]{tex}
	\mathsf{C}_{ijkl}=\boldsymbol{\mathsf{C}}.
\end{minted}

\item 
\begin{equation}
\mathcal{A} \neq \boldsymbol{\mathcal{A}}
\end{equation}
Usando

\begin{minted}[bgcolor=bg]{tex}
	\mathcal{A} \neq \boldsymbol{\mathcal{A}}
\end{minted}


\end{itemize}
\end{itemize}
\item La siguiente expresi\'on
\begin{equation}
 |x| = \left\{ \begin{array}{ll}
         x & \mbox{si $x \geq 0$},\\
        -x & \mbox{si $x < 0$}.\end{array} \right.
\end{equation}

%\item Siga entrenando sus poderes en lenguaje \LaTeX, reproduciendo estas expresiones:
%\begin{equation}
%\Lambda^{\mu'}{}_\nu = \begin{pmatrix}
%\gamma & -\beta\gamma/c & 0 & 0\\
%-\beta\gamma c & \gamma & 0 & 0\\
%0 & 0 & 1 & 0\\
%0 & 0 & 0 & 1 \end{pmatrix},
%\end{equation}
%\begin{equation}
%U^\mu = \frac{\text{d}x^\mu}{\text{d}\tau} =
%\begin{pmatrix} \gamma \\ \gamma v_x \\ \gamma v_y \\ \gamma v_z \end{pmatrix}.
%\end{equation}

\item En su archivo de trabajo, agregue algunas referencias bibliogr'aficas, siguiendo el ejemplo en la página página 53 de la \href{https://udec.instructure.com/courses/40179/pages/latex-pdf-presentacion?module_item_id=1465499}{presentación de clases}.

\item Agregue el paquete de idiomas \texttt{babel}, usando \verb|\usepackage[spanish]{babel}| y vea c'omo afecta al resultado obtenido.

\item Agregue ahora el paquete \texttt{hyperref} agregando 
\verb|\usepackage[colorlinks]{hyperref}| a su archivo (antes del comando \verb|\begin{document}|). Este paquete agrega autom'aticamente \textit{hyperlinks} a su pdf. Verique que ahora puede hacer click en los n'umeros de ecuaciones citados, as'i como en los n'umeros correspondientes a las referencias bibliogr'aficas. Note que \texttt{hyperref} tambi'en agrega hyperlinks a la tabla de contenidos de su pdf, si 'esta existe.

\item El paquete \texttt{hyperref} tambi'en permite incluir \textit{hyperlinks externos} a su pdf. Existen b'asicamente dos formas: La m'as simple es incluir comandos como \verb|\url{http://www.cfm.cl}|, que agrega un hyperlink a la direcci'on se\~nalada. Agregue un ejemplo de este tipo a su archivo de trabajo y vea el resultado.

\item La segunda forma de agregar hyperlinks usando \texttt{hyperref} es con un comando de la forma \verb|\url{link}{texto}|, por ejemplo \verb|\href{http://www.cfm.cl}{FCFM}|, que genera un link a la misma direcci'on anterior, pero que ahora aparece bajo el texto ``FCFM'' en el pdf. Incluya un ejemplo de esto en su archivo de trabajo.

\item Vea el video resumen sobre la creación de tablas en \LaTeX, disponible en Canvas \href{https://udec.instructure.com/courses/40179/pages/tablas?module_item_id=1465510}{aquí}. Siga las sugerencias ahí descritas para familiarizarse con las distintas posibilidades disponibles.

\item Vea el video resumen sobre cómo incluir figuras/gráficos en un documento \LaTeX, disponible en Canvas \href{https://udec.instructure.com/courses/40179/pages/figuras?module_item_id=1465511}{aquí}.
\item En alg'un archivo .tex que ya tenga hecho, inserte el escudo de la UdeC, contenido en el archivo \url{http://www.udec.cl/normasgraficas/sites/default/files/marcaderecha.png}:
\begin{figure}[h!]
\begin{center}
\includegraphics[width=4cm]{marcaderecha.png}
\end{center}
\caption{Marca alineaci'on derecha, Formato PNG.}
\label{fig:escudo}
\end{figure}

Vea \href{http://www.udec.cl/normasgraficas/node/7}{esta} p'agina para conocer otras variaciones oficiales del escudo de la UdeC. 
Puede encontrar m'as informaci'on acerca de las normas gr'aficas de nuestra Universidad en \href{http://www.udec.cl/normasgraficas}{esta} p'agina.

\item Ahora su misi'on es escribir su primer ``art'iculo cient'ifico'' en \LaTeX, que debe reproducir lo m'as fielmente posible \href{https://github.com/gfrubi/CC/blob/master/guias/06/ejemplo-articulo.pdf}{este} ejemplo. Para eso, use todo lo aprendido (en particular, use referencias cruzadas a las ecuaciones, tablas, figuras y referencias). El archivo .pdf de la figura lo puede descargar desde \href{https://github.com/gfrubi/CC/blob/master/guias/06/fig-ajuste-lineal.pdf}{aqu\'i}.

El t'itulo, autor(a), y resumen puede ser incorporados usando los siguientes comandos despu'es del conocido \verb|\documentclass|:

\begin{minted}[bgcolor=bg]{tex}
\title{Escribiendo mi primer artículo con formato Científico, en \LaTeX}
\author{Su nombre (autor(a))}

\begin{document}

\maketitle
\begin{abstract}
Este es el resumen del artículo...
\end{abstract}
\end{minted}
\end{questions}
\end{document} 